\exercisesection{Invariant Subspaces}

\begin{enumerate}
\def\labelenumi{\arabic{enumi})}
\tightlist
\item
  Let \((a_{n})_{n\in \mathbf{N}}\) be a sequence of complex numbers such that there exists a real number \(A\geqslant 0\) satisfying \(n|a_n|\leqslant A\) for all \(n\in \mathbf{N}\) . For \(z\in \mathbf{C}\) of modulus \(<  1\) , set \(\varphi (z) = \sum_{n\geqslant 0}a_nz^n\) .
\end{enumerate}

One assumes that there exists \(\ell \in \mathbf{C}\) such that

\[
\lim_{\substack{x\in \mathbf{R}_{+}\\ x <   1}}\varphi (x) = \ell .
\]

For \(t\in \mathbf{R}_+^*\) , set

\[
g (t) = \sum_ {0 \leqslant n \leqslant t ^ {- 1}} a _ {n}, \quad \mathrm{F} (t) = \varphi (e ^ {- t}).
\]

We denote by \(dx\) the Lebesgue measure on \(\mathbf{R}\) as well as its restriction to \(\mathbf{R}_{+}^{*}\) .

\begin{enumerate}
\def\labelenumi{\alph{enumi})}
\tightlist
\item
  Let \(s \leqslant t\) be elements of \(\mathbf{R}_{+}^{*}\) . We have
\end{enumerate}

\[
| g (s) - g (t) | \leqslant | a _ {[ t ^ {- 1} + 1 ]} | + A \log \left(\frac {t}{s}\right).
\]

\begin{enumerate}
\def\labelenumi{\alph{enumi})}
\setcounter{enumi}{1}
\item
  The function g is slowly oscillating for \(t \rightarrow 0\) .
\item
  For every \(t \in R_{+}^{*}\) , the function \(x \mapsto g(t/x)xe^{-x}\) is integrable on \(R_{+}^{*}\) with respect to the Haar measure \(x^{-1}dx\) .
\item
  For \(0 < x < 1\) , one has
\end{enumerate}

\[
\varphi (x) = \sum_ {n \geqslant 0} g \left(\frac {1}{n}\right) \left(x ^ {n} - x ^ {n + 1}\right).
\]

\begin{enumerate}
\def\labelenumi{\alph{enumi})}
\setcounter{enumi}{4}
\tightlist
\item
  Let \(f(x) = xe^{-x}\) for \(x \in R_{+}^{*}\) . We have \(f \in \mathrm{L}^{1}(\mathbf{R}_{+}^{*}, dx/x)\) ; its integral is equal to 1, and we have F = g * f.
\end{enumerate}

\(f)\) The function \(g\) belongs to \(\mathrm{L}^{\infty}(\mathbf{R}_{+}^{*})\) . (Show that \(\mathrm{F}-g\) is bounded using \(a\) ) as well as the hypothesis.)

\(g)\) The Fourier transform of the function \(f\) on the group \(\mathbf{R}_{+}^{*}\) does not vanish. (Identify \(\mathcal{F}(f)\) to the function \(y\mapsto\Gamma(1+iy)\) and use FVR, VII.)

\begin{enumerate}
\def\labelenumi{\alph{enumi})}
\setcounter{enumi}{7}
\tightlist
\item
  We have
\end{enumerate}

\[
\sum_ {n = 0} ^ {+ \infty} a _ {n} = \ell .
\]

This statement is Littlewood's "Tauberian" theorem. The hypothesis concerning the sequence \((a_n)\) , that is, \(a_n = \mathrm{O}(1/n)\) , can moreover be weakened: it suffices that there exist a constant \(A \geqslant 0\) such that \(na_n \geqslant -A\) for every integer \(n \geqslant 1\) (“Hardy–Littlewood Tauberian theorem,” see for example J. KOREVAAR, Tauberian theory, Grundlehren math. wiss. 329, Springer (2004), Th. 7.4, p.~16; the reader should take care not to confuse this theorem with that of FVR, I, p.~50, Exercise 18).

\begin{enumerate}
\def\labelenumi{\arabic{enumi})}
\setcounter{enumi}{1}
\item
  Let F be a closed subset of \(\widehat{\mathbf{G}}\) . Let K be a compact subset of \(\widehat{\mathbf{G}}\) such that \(\mathrm{F} \cap \mathrm{K} = \varnothing\) . Show that there exists a continuous function \(g\) integrable on G such that \(\mathcal{F}(g)\) is equal to 0 on F and to 1 on K. (Let \(f \in \mathrm{L}^{1}(\mathrm{G})\) be such that \(\mathcal{F}(f)\) is zero on F and equal to 1 on K; consider \(f * \varphi\) , where \(\varphi\) is the Fourier transform of a function of the form \(h * h'\) with \(h \in \mathcal{K}(\widehat{\mathrm{G}})\) , \(h' \in \mathcal{K}(\widehat{\mathrm{G}})\) , and where \(h * h' = 1\) on K.)
\item
  Let \(g\) be a complex function defined on \(]0, +\infty[\) , integrable with respect to Lebesgue measure. We assume that \(\int_0^{+\infty} g(t)dt = 1\) and that \(\int_0^{+\infty} g(t)t^{ix}dt \neq 0\) for all \(x \in \mathbf{R}\) . Let \(f\) be a complex measurable and bounded function in \(]0, +\infty[\) , such that \(\frac{1}{x}\int_0^{+\infty} g(t/x)f(t)dt\) tends to a finite limit \(\ell\) when \(x\) tends to \(+\infty\) .
\end{enumerate}

\begin{enumerate}
\def\labelenumi{\alph{enumi})}
\tightlist
\item
  For any complex-valued function h defined on \(]0, +\infty[\) , integrable with respect to Lebesgue measure and with integral 1, we have
\end{enumerate}

\[
\lim _ {x \to + \infty} \frac {1}{x} \int_ {0} ^ {+ \infty} h \left(\frac {t}{x}\right) f (t) d t = \ell .
\]

In particular,

\[
\lim _ {x \rightarrow + \infty} \frac {1}{x} \int_ {0} ^ {x} f (t) d t = \ell .
\]

\((\text{Setting } g_1(t) = tg(t), h_1(t) = th(t), \text{ we have } g_1 \in \mathrm{L}^1(\mathbf{R}_+^*) ; \text{ the Fourier transform } \mathcal{F}(g_1) \text{ does not vanish and } \check{g}_1 * f \text{ tends to } \ell, \text{ hence } h_1 * f \text{ tends to } \ell.)\) b) If \(f(t)\) is slowly oscillating on \(\mathbf{R}_+^*\) when \(t\) tends to \(+\infty\) , then \(f\) tends to \(\ell\) when \(t\) tends to \(+\infty\) .

\begin{enumerate}
\def\labelenumi{\arabic{enumi})}
\setcounter{enumi}{3}
\tightlist
\item
  Let G = R. For every real number \(\alpha > 1\) , we define \(f_{\alpha} \in L^{1}(\mathbf{R})\) by
\end{enumerate}

\[
f _ {\alpha} (x) = \left\{ \begin{array}{l l} 2 & \text { if } 0 <   x <   1 \\ 1 & \text { if } 1 \leqslant x <   \alpha \\ 0 & \text { otherwise. } \end{array} \right.
\]

\begin{enumerate}
\def\labelenumi{\alph{enumi})}
\item
  The functions \(f * \varepsilon_{x}\) for \(x \in R\) form a total set in \(\mathrm{L}^{1}(\mathbf{R})\) if and only if \(\alpha\) is irrational.
\item
  The functions \(f * \varepsilon_x\) for \(x \in \mathbf{R}\) form a total set in \(L^2(\mathbf{R})\) .
\end{enumerate}

\begin{enumerate}
\def\labelenumi{\arabic{enumi})}
\setcounter{enumi}{4}
\tightlist
\item
  Let \(f \in \mathrm{L}^{\infty}([0, +\infty([)\text{ and } k \in \mathbf{R}_{+}^{*}. \text{ If}\)
\end{enumerate}

\[
\lim _ {x \to + \infty} \frac {1}{x} \int_ {0} ^ {+ \infty} e ^ {- t / x} f (t) d t = \ell
\]

then

\[
\lim _ {x \rightarrow + \infty} \frac {k}{x ^ {k}} \int_ {0} ^ {x} (x - t) ^ {k - 1} f (t) d t = \ell
\]

(Use the functions \(g(t) = e^{-t}\) , and \(h(t) = k(1 - t)^{k-1}\) for \(0 \leqslant t \leqslant 1\) , \(h(t) = 0\) for \(t > 1\) .)

\begin{enumerate}
\def\labelenumi{\arabic{enumi})}
\setcounter{enumi}{5}
\tightlist
\item
  Let \(f \colon \mathbf{R} \to \mathbf{R}\) and \(g: \mathbf{R} \to \mathbf{C}\) be two functions. We assume that \(f \in \mathrm{L}^{\infty}(\mathbf{R})\) and that \(g \in \mathrm{L}^{1}(\mathbf{R})\) , \(\int_{\mathbf{R}} g(x) dx = 1\) and that \(\mathcal{F}(g)\) does not vanish. We assume that \(f\) is slowly decreasing, that is, for \(y > x\) such that \(y - x\) tends to 0 and \(x \to +\infty\) , one has
\end{enumerate}

\[
\liminf (f (y) - f (x)) \geqslant 0.
\]

If \((g*f)(x)\) tends to \(\ell\) as x tends to \(+\infty\) , then \(f(x)\) tends to \(\ell\) as x tends to \(+\infty\) .

\begin{enumerate}
\def\labelenumi{\arabic{enumi})}
\setcounter{enumi}{6}
\tightlist
\item
  Let \(g: R \rightarrow C\) be a continuous function such that:
\end{enumerate}

\[
\sum_ {n = - \infty} ^ {+ \infty} \sup _ {n \leqslant t \leqslant n + 1} | g (t) | <   + \infty .
\]

\begin{enumerate}
\def\labelenumi{\alph{enumi})}
\item
  We have \(g \in \mathrm{L}^{1}(\mathbf{R})\) .
\item
  Suppose that g has integral 1 and that \(\mathcal{F}(g)\) does not vanish. Let \(\mu\) be a complex measure on R such that \(|\mu|([x,x+1])\) is bounded as x ranges over R. If \((g*\mu)(x)\) tends to \(\ell\) as x tends to \(+\infty\) then, for every continuous function \(h:R\to C\) such that
\end{enumerate}

\[
\sum_ {n = - \infty} ^ {+ \infty} \sup _ {n \leqslant t \leqslant n + 1} | h (t) | <   + \infty ,
\]

of integral 1, the limit of \((h * \mu)(x)\) when \(x\) tends to \(+\infty\) is equal to \(\ell\) . (Show that \(h * \mu\) is slowly oscillating and that \((g * (h * \mu))(x)\) tends to \(\ell\) when \(x\) tends to \(+\infty\) .)

\begin{enumerate}
\def\labelenumi{\arabic{enumi})}
\setcounter{enumi}{7}
\tightlist
\item
  Let \(g\) be a function satisfying the hypotheses of Exerc. 3. Moreover, suppose \(g \geqslant 0\) and \(\liminf_{x \to 0} g(x) > 0\) . Let \(f\) be a real-valued measurable function bounded below on \(]0, +\infty[\) . If
\end{enumerate}

\[
\lim _ {x \to + \infty} \frac {1}{x} \int_ {0} ^ {+ \infty} g \left(\frac {t}{x}\right) f (t) d t = \ell
\]

then

\[
\sigma (x) = \frac {1}{x} \int_ {0} ^ {x} f (t) d t
\]

tends to \(\ell\) when \(x\) tends to \(+\infty\) . (Show that \(\sigma\) is bounded for \(x \geqslant 1\) , then that \(\sigma\) is slowly decreasing, then that

\[
\frac {1}{x} \int_ {0} ^ {+ \infty} g \left(\frac {t}{x}\right) \sigma (t) d t
\]

tends to \(\ell\) when \(x\) tends to \(+\infty\) .)

\begin{enumerate}
\def\labelenumi{\arabic{enumi})}
\setcounter{enumi}{8}
\tightlist
\item
  Let \(\varphi\colon[0,+\infty[\to\mathbf{R}\) be a nonnegative increasing function, such that the function \(t\mapsto e^{-\sigma t}\varphi(t)\) is integrable with respect to Lebesgue measure when \(\sigma>1\) . For every \(s\in\mathbf{C}\) such that \(\mathcal{R}(s)>1\) , set
\end{enumerate}

\[
f (s) = \int_ {0} ^ {+ \infty} e ^ {- s t} \varphi (t) d t
\]

(“Laplace transform of \(\varphi\)”).

Suppose there exists \(A \in C\) with the following property: as \(\sigma\) tends to 1 through values \textgreater{} 1, the function

\[
t \mapsto f (\sigma + i \tau) - \frac {\mathrm{A}}{\sigma + i \tau - 1} \qquad \qquad \tau \in \mathbf {R}
\]

converges uniformly to a function g on every compact subset of R. Then

\[
\lim _ {t \to + \infty} \varphi (t) e ^ {- t} = A
\]

(“Ikehara tauberian theorem”).

(We set \(a(t) = e^{-t}\varphi(t)\) for \(t > 0\) , \(a(t) = 0\) for \(t \leqslant 0\) , \(A(t) = A\) for \(t > 0\) , \(A(t) = 0\) for \(t \leqslant 0\) . Let \(\lambda > 0\). We set:

\[
k _ {\lambda} (t) = \lambda \sqrt {\frac {2}{\pi}} \Big (\frac {\sin 2 \lambda t}{2 \lambda t} \Big) ^ {2}
\]

\[
\mathrm{K} _ {\lambda} (t) = \left\{ \begin{array}{l l} 1 - \left| \frac {t}{2 \lambda} \right| & \text { if } | t | \leqslant 2 \lambda \\ 0 & \text { if } | t | > 2 \lambda . \end{array} \right.
\]

Using exercise 1 of II, p.~262, the Lebesgue–Fubini theorem, and the Plancherel theorem, show that:

\[
\lim _ {\varepsilon \to 0} \frac {1}{\sqrt {2 \pi}} \int_ {- \infty} ^ {+ \infty} k _ {\lambda} (x - t) (a (t) - \mathrm{A} (t)) e ^ {- \varepsilon t} d t = \frac {1}{\sqrt {2 \pi}} \int_ {- 2 \lambda} ^ {2 \lambda} \mathrm{K} _ {\lambda} (y) e ^ {- i x y} g (y) d y
\]

and deduce that

\[
\lim _ {x \to + \infty} \frac {1}{\sqrt {2 \pi}} \int_ {- \infty} ^ {+ \infty} k _ {\lambda} (x - t) a (t) d t = \mathrm{A}.
\]

Prove then that \(a\) is bounded and slowly decreasing. One can finally apply exerc. 6.)

\begin{enumerate}
\def\labelenumi{\arabic{enumi})}
\setcounter{enumi}{9}
\tightlist
\item
  We define the von Mangoldt function on \(\mathbf{N}^{*}\) by \(\Lambda(n)=0\) if n is not of the form \(p^{k}\) where p is prime and \(k\geqslant1\) , and \(\Lambda(p^{k})=\log(p)\) .
\end{enumerate}

We denote by \(\zeta\) the holomorphic function on \(\mathbf{C} - \{1\}\) whose restriction to the set of \(s \in \mathbf{C}\) such that \(\mathcal{R}(s) > 1\) coincides with the Riemann zeta function (Exercise 16 of II, p.~266).

For \(x > 0\), we set

\[
\psi (x) = \sum_ {1 \leqslant n \leqslant x} \Lambda (n).
\]

\begin{enumerate}
\def\labelenumi{\alph{enumi})}
\tightlist
\item
  Show that for all \(s \in \mathbf{C}\) such that \(\mathcal{R}(s) > 1\) , one has
\end{enumerate}

\[
\sum_ {n \geqslant 1} \Lambda (n) n ^ {- s} = - \frac {\zeta^ {\prime} (s)}{\zeta (s)}.
\]

\begin{enumerate}
\def\labelenumi{\alph{enumi})}
\setcounter{enumi}{1}
\tightlist
\item
  We have
\end{enumerate}

\[
\lim _ {x \to + \infty} \frac {\psi (x)}{x} = 1.
\]

(Use Exercise 16 of II, p.~266 and Exercise 9).

\begin{enumerate}
\def\labelenumi{\alph{enumi})}
\setcounter{enumi}{2}
\tightlist
\item
  Let \(\pi(x)\) be the number of prime numbers \(\leqslant x\) . Show that
\end{enumerate}

\[
\pi (x) \sim \frac {x}{\log (x)}
\]

when \(x \to +\infty\) (“prime number theorem,” due to Hadamard and de la Vallée Poussin).

\begin{enumerate}
\def\labelenumi{\arabic{enumi})}
\setcounter{enumi}{10}
\tightlist
\item
  One defines the Möbius function \(\mu\) on \(\mathbf{N}^*\) by \(\mu(n) = (-1)^k\) if \(n\) is the product of \(k\) prime factors pairwise distinct and \(\mu(n) = 0\) otherwise.
\end{enumerate}

For \(x > 0\), we set

\[
\mathrm{M} (x) = \sum_ {n \leqslant x} \mu (n), \quad f (x) = \frac {\mathrm{M} (x)}{x}.
\]

\begin{enumerate}
\def\labelenumi{\alph{enumi})}
\tightlist
\item
  Show that for \(\mathcal{R}(s) > 1\) , one has
\end{enumerate}

\[
\sum_ {n \geqslant 1} \mu (n) n ^ {- s} = \frac {1}{\zeta (s)}.
\]

\begin{enumerate}
\def\labelenumi{\alph{enumi})}
\setcounter{enumi}{1}
\item
  Show that f is bounded, and that \(f(y) - f(x) = \mathrm{O}((y - x)/x)\) for \(y \geqslant x \geqslant 1\) , so that f is slowly oscillating on the group \(R_{+}^{*}\) as x tends to \(+\infty\) .
\item
  Let \(a, b \in R_{+}^{*}\) . For \(x \textgreater{} 0\), let {[}x{]} be the integer part of x. We set \(g_{0}(t) = [t^{-1}]\) for \(t \in R_{+}^{*}\) and
\end{enumerate}

\[
g (t) = 2 g _ {0} (t) - a g _ {0} (a t) - b g _ {0} (b t).
\]

The function g is integrable on \(R_{+}^{*}\) with respect to Lebesgue measure. If \(s \in C\) and \(\mathcal{R}(s) > 0\) , one has

\[
\int_ {0} ^ {+ \infty} g (t) t ^ {s} d t = (2 - a ^ {- s} - b ^ {- s}) \frac {\zeta (1 + s)}{1 + s}.
\]

\(d)\) We have \(\int_{0}^{+\infty} g(t)t^{ix} dt \neq 0\) for all \(x \in \mathbf{R}\) . (Use the part \(j\) ) of Exercise 16 of II, p.~266.)

\begin{enumerate}
\def\labelenumi{\alph{enumi})}
\setcounter{enumi}{4}
\tightlist
\item
  Show that \(\int_0^{+\infty} g(t / x)f(t)dt = o(x)\) when \(x \to +\infty\) . Deduce that \(\mathrm{M}(x) = o(x)\) when \(x \to +\infty\) .
\end{enumerate}

(One can show in an elementary way that this result is equivalent to the prime number theorem, cf.~Exercise 10; see for example H. IWANIEC and E. KOWALSKI, Analytic number theory, A.M.S Coll. Publ. 53 (2004), pp.~31--32.)

\begin{enumerate}
\def\labelenumi{\arabic{enumi})}
\setcounter{enumi}{11}
\item
  Let J be the set of \(f \in \mathcal{S}(\mathbf{R}^3)\) such that \(\mathcal{F}(f)\) vanishes on the sphere \(\mathbf{S}_2\) . Let I be the set of \(f \in \mathrm{J}\) such that \((\partial/\partial y_1)(\mathcal{F}(f))\) vanishes on \(\mathbf{S}_2\) . Let \(\overline{\mathrm{I}}\) and \(\overline{\mathrm{J}}\) be the closures of I and J in \(\mathrm{L}^1(\mathbf{R}^3)\) . Then we have \(\mathrm{V}(\overline{\mathrm{I}}) = \mathrm{V}(\overline{\mathrm{J}}) = \mathbf{S}_2\) but \(\overline{\mathrm{I}} \neq \overline{\mathrm{J}}\) . (Let \(\mu\) be the positive measure of mass 1 on \(\mathbf{S}_2\) invariant under the orthogonal group of \(\mathbf{R}^3\) ; show, using Exercise 20 of II, p.~269, c), that \(f(t) = t_1\mathcal{F}(\mu)(t)\) is an element of \(\mathrm{L}^\infty(\mathbf{R}^3)\) orthogonal to I but not to J.)
\item
  One calls a C-set in \(\widehat{\mathbf{G}}\) a closed subset E of \(\widehat{\mathbf{G}}\) having the following property: if \(f\in \mathrm{L}^{1}(\mathrm{G})\) , if \(\mathcal{F}(f)\) vanishes on E, and if \(\varepsilon >0\) , there exists a function \(g\in \mathrm{L}^{1}(\mathrm{G})\) such that \(\| f - f*g\| \leqslant \varepsilon\) and such that \(\operatorname {Supp}(\mathcal{F}(g))\) be compact and disjoint from E.
\end{enumerate}

\begin{enumerate}
\def\labelenumi{\alph{enumi})}
\item
  If E is a C-set in \(\widehat{G}\) , there exists only one closed ideal I of \(\mathrm{L}^1 (\mathrm{G})\) such that \(\mathrm{V(I)} = \mathrm{E}\) .
\item
  Every singleton subset of \(\widehat{\mathrm{G}}\) is a C-set.
\item
  The union of two C-sets is a C-set.
\item
  Every set obtained by translation from a C-set is a C-set.
\item
  Let H be a closed subgroup of \(\widehat{\mathbf{G}}\) , let E be a closed subset of H, and let E' be the boundary of E relative to H. If E' is a C-set relative to \(\widehat{\mathbf{G}}\) , then E is a C-set with respect to \(\widehat{\mathbf{G}}\) . (Use exerc. 21 of II, p.~270.)
\item
  In \(R^{n}\) , every intersection E of a finite number of closed half-spaces is a C-set. (Reason by induction on the dimension of the affine subspace generated by E, using b), c), d), e).
\end{enumerate}

\begin{enumerate}
\def\labelenumi{\arabic{enumi})}
\setcounter{enumi}{13}
\tightlist
\item
  Let E be a closed subset of \(R^{n}\) . Suppose that there exists an interior point p of E such that every line passing through p meets the boundary of E in at most two points. Then there exists only one closed ideal I of \(L^{1}(\mathbf{R}^{n})\) such that \(\mathcal{F}(f)\) vanishes on E. (Consider f as a limit in \(L^{1}(\mathbf{R}^{n})\) of functions deduced from f by homotheties with center p).
\end{enumerate}

¶ 15) We denote by μ a Haar measure of G and by B the filter of neighborhoods of e in \(\widehat{G}\) .

\begin{enumerate}
\def\labelenumi{\alph{enumi})}
\tightlist
\item
  Let U be an integrable open neighborhood of e in G. There exists a function \(a \in \mathrm{L}^{1}(\mathrm{G})\) , \(a \geqslant 0\) , of integral 1, vanishing outside U, such that \(\mathcal{F}(a) \in \mathrm{L}^{1}(\widehat{\mathrm{G}})\) , and
\end{enumerate}

\[
\int_ {\mathrm{G}} a (x) ^ {2} d x \leqslant \frac {2}{\mu (\mathrm{U})}.
\]

(Let V be a compact neighborhood of e such that \(V \subset U\) , \(\mu(U) < \sqrt{2}\mu(V)\) ; let W be a compact symmetric neighborhood of e such that \(VW^{2} \subset U\) ; take \(a = \lambda f * g\) , where \(\lambda\) is a constant \textgreater{} 0, and where f, g are the characteristic functions of the sets VW and W.)

\begin{enumerate}
\def\labelenumi{\alph{enumi})}
\setcounter{enumi}{1}
\tightlist
\item
  If \(f \in \mathrm{L}^{\infty}(\mathrm{G})\) , we denote by \(\mathrm{A}(f)\) the set of elements of \(\widehat{G}\) which belong to the weakly closed vector subspace of \(\mathrm{L}^{\infty}(\mathrm{G})\) invariant under translation generated by f. If \(f, g \in \mathrm{L}^{\infty}(\mathrm{G})\) , one has
\end{enumerate}

\[
\mathrm{A} (f g) \subset \overline {{\mathrm{A} (f) \mathrm{A} (g)}}.
\]

(Let U be a neighborhood of \(e\) . Show that \(fg\) is the weak limit of linear combinations of elements of \(\widehat{G}\) belonging to \(\mathrm{A}(f)\mathrm{UA}(g)\mathrm{U}\) .)

\begin{enumerate}
\def\labelenumi{\alph{enumi})}
\setcounter{enumi}{2}
\tightlist
\item
  Let \(g \in \mathrm{L}^{\infty}(\mathrm{G})\) , K a compact subset of \(\widehat{G}\) , and f a function in \(\mathrm{L}^{1}(\mathrm{G})\) such that \(\mathcal{F}(f)\) vanishes on \(\mathrm{A} = \mathrm{A}(g)\) and outside K. For every \(V \in B\) , let \(\omega(V)\) be the supremum of \(|\mathcal{F}(f)|\) on AV. If
\end{enumerate}

\[
\operatorname * {l i m i n f} _ {\mathfrak {B}} \frac {\omega (\mathrm{V}) ^ {2} \mu ((\mathrm{AV} - \mathrm{A}) \cap \mathrm{K})}{\mu (\mathrm{V})} = 0,
\]

then \(\langle f,g\rangle=0\) . (Let \(\varepsilon>0\) . There exists a compact subset H of G such that \(\int_{G-H}|f|dx\leqslant\varepsilon\) . Then there exists an open neighborhood U of e in \(\widehat{G}\) such that \(|\langle x,\widehat{x}\rangle-1|\leqslant\varepsilon\) for \(x\in H\) , \(\widehat{x}\in U\) . Applying a) to \(\widehat{G}\) and U, one obtains a function \(a\in L^{1}(\widehat{G})\) . Let \(b=\mathcal{F}(a)\in L^{1}(G)\) . Using b), show that \(\mathcal{F}(bg)\) vanishes outside AU. On the other hand,

\[
\left| \int_ {\mathrm{G}} f g d x - \int_ {\mathrm{G}} f g b d x \right| \leqslant \varepsilon (\| f \| _ {1} + 2) \| g \| _ {\infty}
\]

and, since \(f, gb \in \mathrm{L}^{2}(\mathrm{G})\) ,

\[
\int_ {\mathrm{G}} f g b d x = \int_ {\widehat {\mathrm{G}}} \mathcal {F} (f) \mathcal {F} (g b) d \widehat {x} = \int_ {(\mathrm{AU-A}) \cap \mathrm{K}} \mathcal {F} (f) \mathcal {F} (g b) d \widehat {x}.
\]

Bound the absolute value of this last expression by:

\[
2 \| g \| _ {\infty} \frac {\omega (\mathrm{U}) \sqrt {\mu ((\mathrm{AU} - \mathrm{A}) \cap \mathrm{K})}}{\sqrt {\mu (\mathrm{U})}}
\]

to conclude.)

\(d)\) Let I be a closed ideal of \(\mathrm{L}^1 (\mathrm{G})\) , \(\mathrm{A} = \mathrm{V}(\mathrm{I})\) , \(f\) a function of \(\mathrm{L}^1 (\mathrm{G})\) such that \(\mathcal{F}(f)\) vanishes on A, S the set of points where \(\mathcal{F}(f)\) does not vanish, \(\mathrm{S}'\) the boundary of S, and B the largest perfect set contained in \(\mathrm{S} \cap \mathrm{A}\) . For \(\mathrm{V} \in \mathfrak{B}\) , let \(\omega (\mathrm{V})\) be the supremum of \(|\mathcal{F}(f)|\) on AV. We assume that every point of B has a neighborhood K in \(\widehat{\mathrm{G}}\) such that

\[
\operatorname * {l i m i n f} _ {\mathfrak {B}} \frac {\omega (\mathrm{V}) ^ {2} \mu ((\mathrm{AV-A}) \cap \mathrm{K})}{\mu (\mathrm{V})} = 0.
\]

Then \(f \in I\) . (Use the fact that \(L^1(G)\) satisfies Ditkin's condition, and argue as in Prop. 5 of I, p.~94; let \(G' = \widehat{G} \cup \{\infty\}\) be the Alexandroff compactification of \(\widehat{G}\) and let N be the set of \(\chi \in G'\) such that \(\mathcal{F}(f)\) does not belong to \(\mathcal{F}(I)\) in a neighborhood of \(\chi\) . First show that \(N \subset B \cup \{\infty\}\) . Then, using \(c\) ), show that \(N \subset \{\infty\}\) . Finally, show that \(N = \varnothing\) .)

\begin{enumerate}
\def\labelenumi{\alph{enumi})}
\setcounter{enumi}{4}
\tightlist
\item
  Let I be a closed ideal of \(\mathrm{L}^{1}(\mathbf{R}^{n})\) , A = V(I), and f a function in \(\mathrm{L}^{1}(\mathbf{R}^{n})\) such that \(\mathcal{F}(f)\) vanishes on A. For h \textgreater{} 0, let \(A_{h}\) be the set of points of \(R^{n}\) outside A and located at a distance from A less than h, and let \(\omega(h) = \sup_{x \in \mathrm{A}_{h}} |\mathcal{F}(f)(x)|\) . If
\end{enumerate}

\[
\liminf _ {h \to 0} \frac {\omega (h) ^ {2} \mu (\mathrm{A} _ {h})}{h ^ {n}} = 0
\]

then \(f\in \mathrm{I}\) . (Use \(d\) ).)

\(f)\) Let \(f \in \mathrm{L}^{1}(\mathbf{R})\) , \(\alpha \in]0,1[\) , and suppose \(\int_{\mathbf{R}} |y|^{\alpha}| f(y)| dy < +\infty\) . Then there exists a constant \(k\) such that \(|\mathcal{F}(f)(x+h)-\mathcal{F}(f)(x)| \leqslant kh^{\alpha}\) for all \(x \in \mathbf{R}\) and every \(h \in \mathbf{R}\) .

\(g)\) Let I be a closed ideal of \(\mathrm{L}^1 (\mathbf{R})\) , \(f\) a function of \(\mathrm{L}^1 (\mathbf{R})\) such that \(\mathcal{F}(f)\) vanishes on \(\mathrm{V(I)}\) . If

\[
\int_ {\mathbf {R}} | y ^ {1 / 2} f (y) | d y <   + \infty ,
\]

then we have \(f \in I\) . (Use parts e) and f).

\begin{enumerate}
\def\labelenumi{\arabic{enumi})}
\setcounter{enumi}{15}
\tightlist
\item
  Let \(f \in \mathrm{L}^{\infty}(\mathrm{G})\) . We say that \(f\) satisfies the spectral synthesis condition if it belongs to \(\mathrm{Y}(\mathrm{A}(\mathrm{W}_f))\) , where \(\mathrm{W}_f\) is the intersection of all translation-invariant weakly closed subspaces of \(\mathrm{L}^{\infty}(\mathrm{G})\) which contain \(f\) .
\end{enumerate}

\begin{enumerate}
\def\labelenumi{\alph{enumi})}
\item
  Let \(\mu \in \mathcal{M}^1 (\widehat{\mathrm{G}})\) and \(f = \overline{\mathcal{F}}_{\widehat{\mathrm{G}}}(\mu)\in \mathrm{L}^{\infty}(\mathrm{G})\) . Then \(\mathrm{A}(\mathrm{W}_f)\) is the support of \(\mu\) . Moreover, \(f\) satisfies the spectral synthesis condition.
\item
  If G is discrete, every function \(f \in \mathrm{L}^{2}(\mathrm{G})\) satisfies the spectral synthesis condition.
\end{enumerate}

\begin{enumerate}
\def\labelenumi{\arabic{enumi})}
\setcounter{enumi}{16}
\tightlist
\item
  For every \(p \geqslant 1\) , we identify \(\mathrm{L}^p(\mathbf{R}_+)\) with a closed subspace of \(\mathrm{L}^p(\mathbf{R})\) . Let \(f \in \mathrm{L}^1(\mathbf{R}_+)\) be the function defined by \(f(x) = 2e^{-x}\) for \(x \in \mathbf{R}_+\) , and let \(g \in \mathrm{L}^1(\mathbf{R})\) be the function such that \(g(x) = f(-x)\) .
\end{enumerate}

\begin{enumerate}
\def\labelenumi{\alph{enumi})}
\item
  Show that \(f + g = f * g\) .
\item
  Let A (resp. B) be the subalgebra of \(\mathrm{L}^{1}(\mathbf{R})\) generated by f (resp. by f and g). Show that A is dense in \(\mathrm{L}^{1}(\mathbf{R}_{+})\) . (Let \(h \in \mathrm{L}^{\infty}(\mathbf{R}_{+})\) such that \(\langle h, f^{*n} \rangle = 0\) for every integer \(n \geqslant 1\) , where \(f^{*n}\) is the n-fold iterated convolution of f; show that the Laplace transform
\end{enumerate}

\[
\mathrm{F} (z) = \int_ {0} ^ {+ \infty} e ^ {- x z} h (x) d x,
\]

defined for all \(z \in C\) such that \(\mathcal{R}(z) > 0\) is zero.) Deduce that B is dense in \(\mathrm{L}^{1}(\mathbf{R})\) .

\begin{enumerate}
\def\labelenumi{\alph{enumi})}
\setcounter{enumi}{2}
\item
  Let C be a proper closed subalgebra of \(\mathrm{L}^{1}(\mathbf{R})\) containing A. We have \(1 \in \mathrm{Sp}_{\mathrm{C}}(f)\) . (Show that otherwise, part (a) and the holomorphic functional calculus in C would imply that \(g \in C\) .)
\item
  There exists a character \(\chi\) of C such that \(\chi(f)=1\) ; we have
\end{enumerate}

\[
\chi (\varphi) = \int_ {0} ^ {+ \infty} e ^ {- x} \varphi (x) d x
\]

for every function \(\varphi\in\mathrm{L}^{1}(\mathbf{R}_{+})\) .

\begin{enumerate}
\def\labelenumi{\alph{enumi})}
\setcounter{enumi}{4}
\item
  There exists a measure \(\mu \in \mathcal{M}^1 (\mathbf{R})\) of norm 1 such that \(\chi (\varphi) = \int \widehat{\varphi} d\mu\) for every function \(\varphi \in \mathrm{C}\) .
\item
  The Fourier transform of \(\mu\) is the function \(x \mapsto e^{-|x|}\) .
\item
  Let \(\varphi \in \mathrm{C}\) . For every \(y > 0\) , one has
\end{enumerate}

\[
\int_ {\mathbf {R}} \varphi (x) (e ^ {- | x + y |} - e ^ {- | x | - y}) d x = 0.
\]

(Using the relation \(\chi (\varphi *\psi) = \chi (\varphi)\chi (\psi)\) for \(\psi \in \mathrm{L}^1 (\mathbf{R}_+)\) .)

\begin{enumerate}
\def\labelenumi{\alph{enumi})}
\setcounter{enumi}{7}
\tightlist
\item
  We have \(C = L^{1}(\mathbf{R}_{+})\) . (Thus, the subalgebra \(L^{1}(\mathbf{R}_{+})\) is a maximal closed subalgebra of \(L^{1}(\mathbf{R})\) .)
\end{enumerate}

\begin{enumerate}
\def\labelenumi{\arabic{enumi})}
\setcounter{enumi}{17}
\tightlist
\item
  Let G be a discrete ordered commutative topological group. We write the group law of G additively, and denote its identity by 0. We denote by \(G^{+}\) the set of elements \(x \geqslant 0\) of G.
\end{enumerate}

\begin{enumerate}
\def\labelenumi{\alph{enumi})}
\item
  The space \(L_{+}^{1}\) of functions \(f \in \mathrm{L}^{1}(\mathrm{G})\) vanishing outside of \(G^{+}\) is a closed subalgebra of \(\mathrm{L}^{1}(\mathrm{G})\) .
\item
  We assume that G is not Archimedean (TG, V, p.~16, exercise 1 of §3). The algebra L \(_{+}^{1}\) is not maximal in L \(^{1}\) (G).
\end{enumerate}

Assume henceforth that G is Archimedean. Let A be a closed subalgebra of \(\mathrm{L}^{1}(\mathrm{G})\) containing \(L_{+}^{1}\) . For every \(x \in G\) , let \(\varphi_{x} \in \mathrm{L}^{1}(\mathrm{G})\) be the characteristic function of \(\{x\}\) . It is invertible in \(\mathrm{L}^{1}(\mathrm{G})\) and its inverse is \(\varphi_{-x}\) .

\begin{enumerate}
\def\labelenumi{\alph{enumi})}
\setcounter{enumi}{2}
\tightlist
\item
  There exists \(y \in G_{+}\) such that \(\varphi_{y} \in A\) and \(\varphi_{-y} = \varphi_{y}^{-1} \notin A\) and there exists a character \(\chi \neq 0\) of A such that \(\chi(y) = 0\) .
\end{enumerate}

\(d)\) For every \(x > 0\) in G, we have \(\chi(x) = 0\) . (There exists an integer \(k \geqslant 1\) such that \(-y + kx > 0\) .)

\begin{enumerate}
\def\labelenumi{\alph{enumi})}
\setcounter{enumi}{4}
\item
  There exists a measure \(\mu \in \mathcal{M}^1 (\mathrm{G})\) of norm 1 such that \(\chi\) is the restriction of \(\mu\) to the algebra A viewed as a subspace of \(\mathcal{C}(\mathrm{G})\) . (For \(f\in \mathrm{A}\) , one has \(|\chi (f)|\leqslant \varrho (f)\) and moreover \(\varrho (f) = \| f\|_{\infty}\) .)
\item
  Let \(\widehat{\mu}\colon\mathrm{G}\to\mathbf{C}\) be the inverse Fourier transform of \(\mu\) . We have \(\widehat{\mu}(x)=0\) if \(x>0\) .
\item
  The measure \(\mu\) is the counting measure on G.
\item
  The subalgebra \(L_{+}^{1}\) is maximal in \(L^{1}(G)\) .
\end{enumerate}

\begin{enumerate}
\def\labelenumi{\arabic{enumi})}
\setcounter{enumi}{18}
\tightlist
\item
  For \(f \in \mathrm{L}^1(\mathbf{R})\) , we denote by \(\mathrm{A}_f\) the closed subalgebra of \(\mathrm{L}^1(\mathbf{R})\) generated by \(f\) .
\end{enumerate}

\begin{enumerate}
\def\labelenumi{\alph{enumi})}
\tightlist
\item
  There exists \(f \in \mathrm{L}^{1}(\mathbf{R})\) such that \(A_{f}\) is a proper maximal subalgebra of \(\mathrm{L}^{1}(\mathbf{R})\) .
\end{enumerate}

In what follows, we assume that \(f \in \mathrm{L}^{1}(\mathbf{R})\) is a function such that \(A_{f}\) is a proper maximal subalgebra of \(\mathrm{L}^{1}(\mathbf{R})\) .

\begin{enumerate}
\def\labelenumi{\alph{enumi})}
\setcounter{enumi}{1}
\item
  We denote \(S = \{0\} \cup \widehat{f}(\mathbf{R}) \subset \mathbf{C}\) . Show that the function \(z \mapsto \overline{z}\) is not a uniform limit on S of polynomial functions. (First show that C-S is not connected.)
\item
  The subalgebra of \(\mathrm{L}^{1}(\mathbf{R})\) generated by f and \(\widetilde{f}\) is dense in \(\mathrm{L}^{1}(\mathbf{R})\) .
\item
  The Fourier transform of f is injective and does not vanish on R. (Show that S cannot have two bounded connected components.)
\end{enumerate}

\backmatter
